\section{Method}
\label{sec:method}

\subsection{Problem Formulation}
\label{sec:method_formulation}

\noindent\textbf{AnyDepth-YOLO detector.}
We build on an AnyDepth-YOLO~\cite{kang2026anydepth}, a YOLOv12-style detector that
is trained so that its depth can be reduced at inference while reusing the same
weights. Each stage---in \emph{both} the backbone and the detection neck---is split
into an \emph{essential} path that always runs and a \emph{skippable} refinement
path that can be dropped; skipping a stage's refinement path therefore shortens the
network end to end, not just in the backbone. This exposes, from a single set of
parameters and with no extra storage, a family of execution depths. Following the
detector's design we use its two extreme configurations: a \emph{super} path that
executes every refinement path (highest accuracy, most compute) and a \emph{base}
path that skips all of them, keeping only the essential stages (shallowest,
cheapest). Concretely, if the network has $N$ skippable refinement paths across the
backbone and neck, the super path uses the skip mask $[\textsc{false}]^N$ and the
base path $[\textsc{true}]^N$; the detection head is shared and unchanged.
The detector is trained once and then kept \emph{frozen} in our work---we never
modify its weights---so the only inference-time decision is which depth to run on
each frame.

\begin{figure*}[!ht]
\centering
\includegraphics[width=0.75\textwidth]{fig_train.pdf}
\caption{Training pipeline. The AnyDepth-YOLO detector is frozen; both paths are
run once per image to obtain the true advantage $A=L_{\text{base}}-L_{\text{super}}$,
and only the router is updated. The router predicts $\hat{A}$ from the multi-scale
backbone features and is trained with $\text{MSE}(\hat{A},A)$ (blue: forward,
red: backward).}
\label{fig:train}
\end{figure*}

\medskip
\noindent\textbf{Routing problem.}
We thus view the detector as having two execution paths: a deep \emph{super}
configuration $\delta_{\text{sup}}$ and a shallow \emph{base} configuration
$\delta_{\text{base}}$. For each input image, we define the \emph{super
advantage} as the reduction in detection loss obtained by using the deep path:
\begin{equation}
A = L_{\text{base}} - L_{\text{super}},
\label{eq:advantage}
\end{equation}
where $L_{\text{base}}$ and $L_{\text{super}}$ are the per-image detection
losses of the two paths. At deployment, the predicted advantage $\hat{A}$
selects the execution configuration by a single threshold $\tau$:
\begin{equation}
\text{config} =
\begin{cases}
\delta_{\text{super}}, & \hat{A} > \tau,\\[2pt]
\delta_{\text{base}},  & \text{otherwise.}
\end{cases}
\label{eq:route}
\end{equation}


\subsection{Router Training}
\label{sec:method_training}
The router is trained entirely offline against the frozen detector
(Fig.~\ref{fig:train}). We first cache,
once, the features and the true advantage of every training image
(Sec.~\ref{sec:method_caching}), then fit a lightweight head
(Sec.~\ref{sec:method_router}) to regress that advantage
(Sec.~\ref{sec:method_loss}); the detector is never updated.

\subsubsection{Offline Feature and Loss Caching}
\label{sec:method_caching}
The detector is kept \emph{frozen} throughout: the router never changes the
backbone, so for every training image the advantage in Eq.~\eqref{eq:advantage} is
a fixed scalar that can be precomputed once. We exploit this by caching, offline,
everything the router needs, which both decouples router training from the
detector and makes each training epoch a cheap pass over precomputed tensors (no
detector forward, so trade-off and hyperparameter sweeps are fast).
 
Concretely, for every image and for \emph{each} of the two paths
($p \in \{\text{base},\text{super}\}$) we run the frozen detector once and
store: (i) multi-scale backbone features and (ii) the per-image detection loss.
The features are tapped at the three backbone pyramid levels that feed the
detection neck (the P3/P4/P5 stages, denoted $\{S_2, S_3, S_4\}$). For the GAP-MLP
router we apply global average pooling to each level and concatenate, giving a
single 768-d descriptor per image and path; for the TinyConv router we instead
apply adaptive average pooling to a small $G \times G$ spatial grid (default
$2 \times 2$) per level and concatenate along channels, giving a
$768 \times G \times G$ tensor that preserves coarse spatial layout. We optionally
also tap the neck (prediction) levels for the feature-source ablation. The cached
loss is the standard detector training loss (classification plus box and
distribution-focal regression) evaluated per image against the ground-truth boxes;
it correlates strongly with AP and serves as the supervision target through the
advantage $A = L_{\text{base}} - L_{\text{super}}$. The result is a compact
cache (train/val splits) of per-image, per-path feature tensors and losses on
which the router is trained.


\subsubsection{Router Architecture}
\label{sec:method_router}
The router takes the cached multi-scale features and outputs the scalar advantage
prediction $\hat{A}$. 
How this head consumes the features is a design choice
rather than a contribution of this work; we consider two variants and compare them
as an ablation (Sec.~\ref{sec:abl_router}).
 
\noindent\textbf{GAP-MLP.} The 768-d global-average-pooled descriptor is normalized
by a batch-normalization layer, projected by a linear layer with ReLU, and mapped
to $\hat{A}$ by a small multilayer perceptron head. This is simple but discards
all spatial layout.
 
\noindent\textbf{TinyConv.} To retain ``where in the frame the hard content is,''
the router consumes the $G \times G$ feature grids instead of a pooled vector. Each
feature group is processed by a tiny spatial projection: a 2D batch-normalization
layer that equalizes the differing magnitudes of the multi-scale features before
they are aggregated, a $1 \times 1$ convolution that reduces the concatenated
channels to a small width with ReLU, a $3 \times 3$ depthwise
convolution (stride 1, padding 1) with ReLU that mixes the coarse spatial cells,
and finally global average pooling to a $d$-dimensional vector. The group vectors
are concatenated and mapped to $\hat{A}$ by a two-layer head (hidden width 128).
The depthwise convolution and the small channel width keep the router extremely
cheap (Table~\ref{tab:overhead}) while the 2D normalization makes the multi-scale
aggregation well-conditioned.

 
\begin{figure*}[!ht]
\centering
\includegraphics[width=0.75\textwidth]{fig_eval.pdf}
\caption{Deployment on video. At each frame the detector runs the chosen path and
the router predicts $\hat{A}$ from its backbone features; the decision
$\hat{A}>\tau$ selects the path (\textsc{base} or \textsc{super}) for the next
frame. Routing is causal and recursive across frames $t,t{+}1,\dots,T$, adding only
the lightweight router forward.}
\label{fig:eval}
\end{figure*}


\subsubsection{Training Objective}
\label{sec:method_loss}
We train the router with a single regression objective on the cached advantages
(Fig.~\ref{fig:train}):
\begin{equation}
L_{\text{acc}} = \text{MSE}(\hat{A}, A).
\label{eq:loss}
\end{equation}
Unlike efficiency-aware routing objectives of the form
$\lambda_{\text{acc}}L_{\text{acc}} + \lambda_{\text{flops}}L_{\text{flops}}
(+\,\lambda_{\text{uni}}L_{\text{uni}})$, our objective contains no efficiency
weight $\lambda$ and no anti-collapse term, so the deployment operating point is
not fixed during training. The advantage is optionally standardized (z-scored)
across each mini-batch so that a few high-loss images do not dominate the gradient.
Because at deployment the router reads the features of whichever path ran on the
previous frame (Sec.~\ref{sec:method_deploy}), training must expose it to both
feature distributions: for each image we sample the assumed previous action as
\textsc{super} with probability $p$ (and \textsc{base} otherwise) and feed the
corresponding cached features together with that action's identity. A balanced
$p{=}0.5$ keeps the prediction calibrated regardless of which path executed
upstream; we ablate $p$ in Sec.~\ref{sec:abl_prevaction}.
Because the target $A$ is a fixed property of each image under the frozen detector,
the regression is a well-posed supervised problem rather than a non-stationary
policy-learning one: there is no sampled action, hence no estimator variance and no
risk of collapse, and no efficiency term that would tie the model to one budget.
The router instead learns to \emph{rank} images by how much the deep path helps.
This ranking is what makes a single trained model span the whole trade-off family:
sweeping the decision threshold over the range of predicted advantages moves the
operating point continuously from all-base (largest threshold) to all-super
(smallest threshold).




\subsection{Router Deployment}
\label{sec:method_deploy}
Once trained, the router is deployed as a causal per-frame controller on top of the
frozen detector. We first describe the basic threshold-based routing loop that
serves a \emph{fixed} operating point (Sec.~\ref{sec:method_deploy_route}), and then
an optional online extension that adapts the threshold to track a \emph{time-varying}
latency budget (Sec.~\ref{sec:method_deploy_budget}).

\subsubsection{Threshold-Based Routing}
\label{sec:method_deploy_route}
At deployment the routing rule of Eq.~\eqref{eq:route} is applied with a single
scalar threshold $\tau$.
Because the router was trained only to predict the
advantage, $\tau$ is a free parameter chosen at serving time: raising $\tau$ sends
fewer images to the deep path (cheaper, lower accuracy) and lowering it sends more,
so sweeping $\tau$ produces the continuous Pareto family from one model. A target
compute budget is met simply by picking the corresponding $\tau$; in practice
$\tau$ can be fixed on a held-out set so that no video labels are needed at serving
time.

For video, routing is causal and per-frame (Fig.~\ref{fig:eval}): the path for
frame $t$ is decided from features that are available before the heavy forward,
namely the features produced by the path actually executed on frame $t{-}1$.
Exploiting temporal continuity (adjacent frames have similar difficulty), this lets
the router commit to a path at the start of each frame without any extra detector
pass, so the only added cost is the lightweight router forward itself.
Algorithm~\ref{alg:routing} states this loop: only the chosen path is run, and a
fixed $\tau$ holds the system at one operating point on the Pareto front.



\subsubsection{Online Latency-Budget Control}
\label{sec:method_deploy_budget}
The fixed-$\tau$ loop serves one operating point. When the target latency budget
itself varies over time---e.g.\ under thermal or battery constraints---the same
router supports tracking it online by adapting $\tau$ in closed loop, without any
retraining. Algorithm~\ref{alg:budget} wraps a proportional--integral (PI)
controller around the routing loop: it maintains an exponential moving average
$\bar{L}$ of the realized per-frame latency, forms the error $e=L^\star(t)-\bar{L}$
against the budget, and nudges $\tau$ by $-k_p e - k_i I$. Since lowering $\tau$
routes more frames to the deep path (raising latency) and raising it does the
opposite, the controller drives the realized latency toward $L^\star(t)$ at every
frame. This optional layer reuses Algorithm~\ref{alg:routing} unchanged as its
inner step; we evaluate its tracking behavior in Sec.~\ref{sec:results_budget}.

\begin{algorithm}[t]
\caption{Causal budget-adaptive depth routing (deployment)}
\label{alg:routing}
\small
\begin{tabular}{@{}r@{\hspace{5pt}}l@{}}
\multicolumn{2}{@{}l}{\textbf{Input:} frozen AnyDepth detector $D$ with \textsc{base}/\textsc{super} paths,}\\
\multicolumn{2}{@{}l}{\hphantom{\textbf{Input:} }router $R$, threshold $\tau$, frames $x_1,\dots,x_T$}\\
\multicolumn{2}{@{}l}{\textbf{Output:} streaming per-frame detections $y_t$}\\[1pt]
$1$ & $\mathit{prev}\leftarrow\textsc{base}$ \hfill{\footnotesize\itshape $\triangleright$ first frame defaults to base}\\
$2$ & \textbf{for} $t=1$ \textbf{to} $T$ \textbf{do}\\
$3$ & \quad $(y,f)\leftarrow$ run path $\mathit{prev}$ on $x_t$ \hfill{\footnotesize\itshape $\triangleright$ detections + tapped features}\\
$4$ & \quad $\hat{A}\leftarrow R(f,\mathit{prev})$ \hfill{\footnotesize\itshape $\triangleright$ predicted advantage $L_{\mathrm{base}}{-}L_{\mathrm{super}}$}\\
$5$ & \quad $a\leftarrow\textsc{super}$ \textbf{if} $\hat{A}>\tau$ \textbf{else} \textsc{base} \hfill{\footnotesize\itshape $\triangleright$ Eq.~\eqref{eq:route}}\\
$6$ & \quad \textbf{if} $a\neq\mathit{prev}$ \textbf{then} $(y,\cdot)\leftarrow$ run path $a$ on $x_t$\\
$7$ & \quad \textbf{emit} $y$ as the detection of frame $t$\\
$8$ & \quad $\mathit{prev}\leftarrow a$\\
\end{tabular}
\end{algorithm}

\begin{algorithm}[t]
\caption{PI latency-budget control of the threshold $\tau$}
\label{alg:budget}
\small
\begin{tabular}{@{}r@{\hspace{5pt}}l@{}}
\multicolumn{2}{@{}l}{\textbf{Input:} target $L^\star(t)$, path latencies $\ell_{\mathrm{base}},\ell_{\mathrm{super}},\ell_{\mathrm{rtr}}$,}\\
\multicolumn{2}{@{}l}{\hphantom{\textbf{Input:} }gains $k_p,k_i$, EMA factor $\beta$}\\[1pt]
$1$ & $\tau\leftarrow\tau_0$;\;\, $I\leftarrow0$;\;\, $\bar{L}\leftarrow\tfrac12(\ell_{\mathrm{base}}{+}\ell_{\mathrm{super}})$\\
$2$ & \textbf{for} $t=1$ \textbf{to} $T$ \textbf{do}\\
$3$ & \quad $a_t\leftarrow$ route $x_t$ with threshold $\tau$ \hfill{\footnotesize\itshape $\triangleright$ Alg.~\ref{alg:routing}, lines 3--6}\\
$4$ & \quad $\ell\leftarrow\ell_{a_t}+\ell_{\mathrm{rtr}}$ \hfill{\footnotesize\itshape $\triangleright$ realized per-frame latency}\\
$5$ & \quad $\bar{L}\leftarrow\beta\bar{L}+(1{-}\beta)\,\ell$ \hfill{\footnotesize\itshape $\triangleright$ EMA of realized latency}\\
$6$ & \quad $e\leftarrow L^\star(t)-\bar{L}$;\quad $I\leftarrow I+e$\\
$7$ & \quad $\tau\leftarrow\mathrm{clip}(\tau-k_p\,e-k_i\,I)$ \hfill{\footnotesize\itshape $\triangleright$ lower $\tau\Rightarrow$ more \textsc{super}}\\
\end{tabular}
\end{algorithm}